\documentclass[12pt]{article}
\usepackage{amsfonts,amsmath,amssymb}
\usepackage{microtype}
\usepackage[hidelinks]{hyperref}

\setlength{\emergencystretch}{3em}
\linespread{1.08}
\newtheorem{exercise}{Exercise}[section]
\newcommand{\C}{\mathbb C}
\newcommand{\R}{\mathbb R}
\newcommand{\Z}{\mathbb Z}
\newcommand{\End}{\operatorname{End}}
\newcommand{\Hom}{\operatorname{Hom}}
\newcommand{\tr}{\operatorname{tr}}
\newcommand{\im}{\operatorname{im}}

\hypersetup{
  pdftitle={Coordinate-Free Mathematics Through Experiment},
  pdfauthor={Paul Gustafson}
}

\title{Coordinate-Free Mathematics Through Experiment}
\author{Paul Gustafson}
\date{August 2026}

\begin{document}
\maketitle
\tableofcontents

\section{Interference, states, and measurement}

Why do observed waves interfere, and why do quantum measurements return
discrete outcomes?

\subsection{Superposition}

A \textbf{real vector space} is an abelian group $V$ with scalar
multiplication by $\R$ satisfying the distributive and associative laws.  A
map $T:V\to W$ is \textbf{linear} if
\[
T(av+bw)=aT(v)+bT(w).
\]
The \textbf{direct sum} $V\oplus W$ is the product vector space.  The
\textbf{quotient} $V/N$ identifies $v$ and $w$ when $v-w\in N$.

\begin{exercise}[String vibrations obey superposition]
For a homogeneous string of length $L$, constant tension $T$, density $\rho$,
small transverse displacement, and fixed endpoints, derive
\[
\rho\,\partial_t^2u=T\,\partial_x^2u,
\qquad u(0,t)=u(L,t)=0.
\]
Prove that the solutions form a vector space and that evolution from initial
data is linear.  Derive the normal modes, their frequencies, and conservation
of
\[
E(t)=\frac12\int_0^L
\bigl(\rho|\partial_tu|^2+T|\partial_xu|^2\bigr)\,dx.
\]
State which small-displacement and no-damping assumptions are required.
\end{exercise}

\begin{exercise}[Fourier modes produce interference fringes]
Solve the fixed-end wave equation by expanding initial data in the orthogonal
functions $\sin(n\pi x/L)$.  Prove Parseval's identity for a finite mode sum
and use it to express the conserved energy as a sum of modal energies.  For
two coherent monochromatic amplitudes $\psi_1,psi_2$, derive
\[
|\psi_1+\psi_2|^2
=|\psi_1|^2+|\psi_2|^2
+2\operatorname{Re}(\overline{\psi_1}\psi_2).
\]
Determine when the cross term gives constructive and destructive
interference, and compare the predicted fringe spacing with Young's
two-source interference geometry \cite{young}.
\end{exercise}

The \textbf{tensor product} $V\otimes W$ is a vector space equipped with a
bilinear map $(v,w)\mapsto v\otimes w$ such that every bilinear map
$V\times W\to X$ factors through a unique linear map $V\otimes W\to X$.

\begin{exercise}[Tensor products model composites]
Construct $V\otimes W$ as a quotient of the free vector space on $V\times W$
and prove its universal property.  Prove uniqueness up to unique isomorphism.
Explain why closure under superposition forces $V\otimes W$ to contain
non-product vectors.
\end{exercise}

\subsection{Spectral measurement}

A \textbf{finite-dimensional Hilbert space} is a finite-dimensional complex
vector space $H$ with a positive-definite Hermitian inner product.  The
\textbf{adjoint} $A^*$ satisfies
$\langle Av,w\rangle=\langle v,A^*w\rangle$.  An operator is
\textbf{self-adjoint} if $A=A^*$ and \textbf{normal} if $A^*A=AA^*$.
An \textbf{orthogonal projection} satisfies $P=P^*=P^2$.

\begin{exercise}[Spectral measurement]
Prove that every normal $A\in\End(H)$ has a unique decomposition
\[
A=\sum_{\lambda\in\sigma(A)}\lambda P_\lambda,
\qquad
P_\lambda P_\mu=0\quad(\lambda\ne\mu),
\qquad
\sum_\lambda P_\lambda=1.
\]
For self-adjoint $A$ and a unit state $\psi$, prove that
$\langle\psi,P_\lambda\psi\rangle$ is a probability distribution.  For
$A=\sigma_3$, compare the two predicted outcomes with the beam splitting in
the Stern--Gerlach experiment \cite{stern-gerlach}.
\end{exercise}

\begin{exercise}[Noncommuting observables are uncertain]
For self-adjoint $A,B$ and a unit vector $\psi$, define
\[
\Delta_\psi A
=\|(A-\langle\psi,A\psi\rangle)\psi\|.
\]
Apply Cauchy--Schwarz to prove
\[
\Delta_\psi A\,\Delta_\psi B
\geq\frac12|\langle\psi,[A,B]\psi\rangle|.
\]
Characterize equality and compute the bound for two Pauli observables.  State
why the result constrains repeated preparations rather than simultaneous
classical values assigned to one particle.
\end{exercise}

For $\psi\in H_A\otimes H_B$, define
$T_\psi:H_A^*\to H_B$ by $T_\psi(f)=(f\otimes1)(\psi)$.

\begin{exercise}[Schmidt spectra measure entanglement]
Apply singular-value decomposition to $T_\psi$ and prove
\[
\psi=\sum_{i=1}^r s_i a_i\otimes b_i
\]
for orthonormal families and $s_i>0$.  Prove that $\psi$ is a product vector
exactly when $r=1$, and identify $s_i^2$ as the nonzero eigenvalues of both
reduced density operators.
\end{exercise}

\subsection{Open evolution}

A \textbf{state} on $M_n(\C)$ is a positive matrix $\rho$ with
$\tr\rho=1$.  A linear map $\Phi:M_n\to M_m$ is \textbf{positive} if it
preserves positive matrices and \textbf{completely positive} if
$1_{M_r}\otimes\Phi$ is positive for every $r$.  A \textbf{quantum channel}
is completely positive and trace preserving.  Its \textbf{Choi matrix} is
\[
C_\Phi=\sum_{i,j}E_{ij}\otimes\Phi(E_{ij}).
\]

\begin{exercise}[Ancillas force complete positivity]
Prove that transposition is positive but not $2$-positive by applying partial
transpose to a maximally entangled two-qubit state.  Prove the equivalences
\[
\Phi\text{ completely positive}
\Longleftrightarrow C_\Phi\geq0
\Longleftrightarrow
\Phi(x)=\sum_aK_axK_a^*.
\]
Construct a Stinespring dilation from the Kraus operators and compare the Choi
state reconstruction with quantum process tomography
\cite{choi,stinespring,altepeter}.
\end{exercise}

\begin{exercise}[No channel clones nonorthogonal states]
Suppose a channel maps $\rho\otimes\sigma$ to $\rho\otimes\rho$ for every
pure state $\rho$.  Use contractivity of trace distance, or monotonicity of
fidelity, to obtain a contradiction from two nonorthogonal states
\cite{wootters-zurek}.
\end{exercise}


\section{Molecular spectra and symmetry}

Why do experimentally observed spectral lines organize into symmetry types
and selection rules?

\subsection{Representations and characters}

A \textbf{representation} of a group $G$ on $V$ is a homomorphism
$\rho:G\to\operatorname{GL}(V)$.  A subspace is \textbf{invariant} if it is
preserved by every $\rho(g)$ and \textbf{irreducible} if it has no nontrivial
invariant subspaces.  An \textbf{intertwiner} $T:V\to W$ satisfies
$T\rho_V(g)=\rho_W(g)T$.

\begin{exercise}[Finite symmetry decomposes into sectors]
For finite $G$, average an inner product and a projection over $G$.  Deduce
Maschke's theorem and prove Schur's lemma over $\C$.  If a Hamiltonian commutes
with $G$, prove that it is block diagonal on
\[
V\cong\bigoplus_\alpha M_\alpha\otimes V_\alpha
\]
and identify the symmetry-enforced degeneracy of each block.
\end{exercise}

The \textbf{character} of $V$ is $\chi_V(g)=\tr(\rho(g))$.  For class
functions define
\[
\langle\chi,\psi\rangle_G
=\frac1{|G|}\sum_{g\in G}\overline{\chi(g)}\psi(g).
\]

\begin{exercise}[Characters recover multiplicities]
Prove the orthogonality of irreducible characters and
\[
m_\alpha=\langle\chi_\alpha,\chi_V\rangle_G.
\]
Prove that
\[
P_\alpha=\frac{\dim V_\alpha}{|G|}
\sum_{g\in G}\overline{\chi_\alpha(g)}\rho(g)
\]
projects onto the $V_\alpha$-isotypic component.
\end{exercise}

\subsection{Water vibrations and selection rules}

For a nonlinear molecule with $N$ nuclei, the displacement representation has
dimension $3N$.  Its vibrational representation is obtained by removing the
translation and rotation representations.  In the electric-dipole model a
transition $V_i\to V_f$ through a dipole component $V_\mu$ is allowed only if
\[
\Hom_G(1,V_f^*\otimes V_\mu\otimes V_i)\ne0.
\]

\begin{exercise}[Water has three symmetry-classified vibrations]
Using the equilibrium point group $C_{2v}$, compute the character of the
nine-dimensional displacement representation of water and prove
\[
V_{\mathrm{vib}}\cong2A_1\oplus B_2.
\]
Determine the infrared- and Raman-active modes.  Compare the predicted mode
count and symmetry types with measured water spectra \cite{shimanouchi}.
Distinguish the empirical equilibrium geometry and force constants from the
representation-theoretic conclusions.
\end{exercise}

\begin{exercise}[Symmetry gives electric-dipole selection rules]
Let $V_i,V_f$ be irreducible energy sectors and let the dipole operator
transform as the spatial vector representation $V_\mu$.  Prove that a
nonzero transition amplitude is equivalent to the occurrence of the trivial
representation in $V_f^*\otimes V_\mu\otimes V_i$.  For rotation symmetry,
use the decomposition of $V_j\otimes V_1$ to derive
$\Delta j=0,\pm1$ with the $0\to0$ exception.  Separate this symmetry rule
from parity, radial overlap, and experimental line strength.
\end{exercise}

\subsection{Continuous and projective symmetry}

A \textbf{matrix Lie group} is a closed subgroup of
$\operatorname{GL}(n,\C)$.  Its \textbf{Lie algebra} consists of the matrices
$X$ for which $e^{tX}$ stays in the group, with bracket $[X,Y]=XY-YX$.
For rotation generators use
\[
[J_i,J_j]=i\varepsilon_{ijk}J_k,
\qquad J^2=J_1^2+J_2^2+J_3^2.
\]

\begin{exercise}[Rotation symmetry produces spin]
Prove that $J^2$ commutes with every $J_i$.  Using
$J_\pm=J_1\pm iJ_2$, classify the finite-dimensional unitary irreducible
representations by $j\in\tfrac12\Z_{\ge0}$ and
$m=-j,-j+1,\ldots,j$.  Verify the spin-$1/2$ representation with Pauli
matrices and relate its two outcomes to Stern--Gerlach splitting.
\end{exercise}

\begin{exercise}[Hydrogen's hidden rotational symmetry]
For the Coulomb Hamiltonian, define angular momentum $L=q\times p$ and the
Runge--Lenz vector in the bound-state model.  Prove that their commutators
close to $\mathfrak{so}(4)$ after rescaling by the negative energy.  Use the
irreducible representation dimensions to recover the $n^2$ orbital
degeneracy.  Identify relativistic, spin, and radiative effects omitted by the
model and compare the resulting splitting hierarchy with hydrogen
spectroscopy \cite{hydrogen}.
\end{exercise}

A \textbf{projective representation} satisfies
\[
U_gU_h=\alpha(g,h)U_{gh},
\]
where associativity makes $\alpha$ a $U(1)$-valued $2$-cocycle.  Rephasing
$U_g$ changes $\alpha$ by a coboundary.

\begin{exercise}[Cohomology classifies projective symmetry]
Prove that equivalence classes of factor systems are classified by
$H^2(G,U(1))$.  For $G=\Z/2\times\Z/2$ and
$U_x=\sigma_1$, $U_z=\sigma_3$, compute the nontrivial class and prove that
every commuting edge operator is scalar.  State the symmetry and gap
assumptions required for protection of a spin-$1/2$ edge doublet.
\end{exercise}


\section{Electromagnetic fields and geometry}

Why do Maxwell fields, charge conservation, and classical motion admit
coordinate-independent formulations?

\subsection{Forms and Maxwell equations}

For a smooth manifold $M$, a \textbf{differential $p$-form} is a smooth
section of $\Lambda^pT^*M$.  The exterior derivative
$d:\Omega^p(M)\to\Omega^{p+1}(M)$ satisfies $d^2=0$ and the graded Leibniz
rule.  For an oriented manifold with boundary, Stokes' theorem is
\[
\int_Md\omega=\int_{\partial M}\omega.
\]

A pseudo-Riemannian metric and orientation determine the \textbf{Hodge star}
by
\[
\alpha\wedge\star\beta
=\langle\alpha,\beta\rangle_g\operatorname{vol}_g.
\]
On four-dimensional spacetime, the electromagnetic field is
$F\in\Omega^2(M)$ and the current is $J\in\Omega^3(M)$.

\begin{exercise}[Maxwell equations force charge conservation]
From
\[
dF=0,
\qquad d\star F=J,
\]
prove $dJ=0$ and $\int_{\partial X}J=0$ for every compact four-region $X$.
Recover the coordinate continuity equation and identify it as a chart
expression of the invariant statement.  Derive the vacuum wave equation and
compare its wave speed with electromagnetic-wave measurements
\cite{hertz}.
\end{exercise}

\subsection{Hamiltonian mechanics}

A \textbf{symplectic form} is a closed nondegenerate two-form $\omega$.  A
Hamiltonian $H$ determines $X_H$ by
\[
\iota_{X_H}\omega=dH.
\]
The \textbf{Poisson bracket} is
$\{f,h\}=\omega(X_f,X_h)$.

\begin{exercise}[Hamiltonian flow preserves phase volume]
Using Cartan's formula, prove
\[
\mathcal L_{X_H}\omega=0
\]
and deduce preservation of $\omega^n/n!$.  Derive Hamilton's equations in
canonical coordinates.  Explain why friction cannot be an autonomous
Hamiltonian flow on the same finite-dimensional phase space.
\end{exercise}

\begin{exercise}[Magnetic curvature produces the Lorentz force]
On $T^*Q$, replace the canonical symplectic form by
$\omega_B=\omega_0+q\pi^*B$ for a closed two-form $B$.  Derive the Lorentz
force from the kinetic Hamiltonian and identify gauge change as a change of
potential leaving $B$ fixed.
\end{exercise}

\subsection{Metric curvature}

A \textbf{pseudo-Riemannian metric} is a smooth nondegenerate symmetric
bilinear form on $TM$.  Its \textbf{Levi--Civita connection} is torsion free
and metric compatible.  Its curvature is
\[
R(X,Y)Z
=\nabla_X\nabla_YZ-\nabla_Y\nabla_XZ-\nabla_{[X,Y]}Z.
\]

\begin{exercise}[Curvature produces interferometer strain]
Derive the coordinate-free Koszul formula and prove existence and uniqueness
of the Levi--Civita connection.  For a variation of geodesics with tangent
$T$ and separation $J$, prove
\[
\nabla_T\nabla_TJ=R(T,J)T.
\]
Apply this equation to a weak transverse plane gravitational wave and derive
the opposite fractional arm-length changes measured by an equal-arm
Michelson interferometer.  Compare the predicted differential strain with
GW150914 \cite{ligo} and identify the weak-field and freely falling mirror
assumptions.
\end{exercise}

\subsection{Connections and gauge fields}

A \textbf{vector bundle} $E\to M$ is locally a product $U\times V$ with
linear transition maps.  A \textbf{connection} is a linear map
\[
\nabla:\Omega^0(M,E)\to\Omega^1(M,E),
\qquad
\nabla(fs)=df\otimes s+f\nabla s.
\]
Its \textbf{curvature} is
$F_\nabla=\nabla^2\in\Omega^2(M,\End E)$.
In a local gauge, $F_A=dA+A\wedge A$.

\begin{exercise}[Curvature is gauge covariant]
For $A^g=g^{-1}Ag+g^{-1}dg$, prove
\[
F_{A^g}=g^{-1}F_Ag,
\qquad d_AF_A=0.
\]
Using an invariant inner product, vary
\[
S(A)=\frac12\int_M\langle F_A\wedge\star F_A\rangle
\]
and derive the Yang--Mills equation.  Recover Maxwell theory for $G=U(1)$ and
distinguish the empirical gauge group and matter representations from the
geometric consequences.
\end{exercise}


\section{Observed phase and quantized transport}

How can a field-free region carry an observable phase, and why is Hall
conductance quantized?

\subsection{Aharonov--Bohm holonomy}

The \textbf{de Rham cohomology} is
\[
H^p_{\mathrm{dR}}(M)
=\ker\bigl(d:\Omega^p(M)\to\Omega^{p+1}(M)\bigr)
/\im\bigl(d:\Omega^{p-1}(M)\to\Omega^p(M)\bigr).
\]
For a $U(1)$ connection $A$ and closed loop $\gamma$, define
\[
\operatorname{Hol}_A(\gamma)
=\exp\!\left(\frac{iq}{\hbar}\int_\gamma A\right).
\]

\begin{exercise}[A field-free loop carries phase]
Outside an ideal solenoid let
\[
M=\R^3\setminus\{(0,0,z):z\in\R\},
\qquad
A=\frac{\Phi}{2\pi}d\theta.
\]
Prove that the phase depends only on $[A]\in H^1_{\mathrm{dR}}(M)$ and
$[\gamma]\in H_1(M)$, and equals
$\exp(iq\Phi/\hbar)$ for a once-winding loop.  Compare the predicted fringe
shift with shielded-flux electron interferometry
\cite{aharonov-bohm,tonomura}.
\end{exercise}

\begin{exercise}[Berry phase is eigenbundle holonomy]
For a smooth family of self-adjoint operators with a simple isolated
eigenvalue, prove that the eigenlines form a bundle.  For a local unit
eigenvector $\psi$, set
\[
\mathcal A=i\langle\psi,d\psi\rangle,
\qquad \mathcal F=d\mathcal A.
\]
Prove the gauge-transformation law and identify
$\exp(i\int_\gamma\mathcal A)$ with holonomy.  For
$H(n)=n\cdot\sigma$, $n\in S^2$, compute the phase as half the enclosed solid
angle and compare it with neutron spin-rotation measurements
\cite{berry-neutron}.
\end{exercise}

\subsection{Chern classes and stable bundles}

For a complex bundle $E$ with connection $\nabla$, define the Chern forms by
\[
\det\!\left(1+\frac{iF_\nabla}{2\pi}\right)
=1+c_1(E,\nabla)+c_2(E,\nabla)+\cdots.
\]
The \textbf{Grothendieck group} $K^0(X)$ is generated by vector bundles with
$[E\oplus F]=[E]+[F]$.

\begin{exercise}[Curvature produces stable characteristic classes]
Using the Bianchi identity, prove that the Chern forms are closed.  For a path
of connections, derive the transgression formula and prove that their de Rham
classes are connection independent.  Prove that
\[
\operatorname{ch}:K^0(X)\longrightarrow H^{\mathrm{even}}_{\mathrm{dR}}(X)
\]
is well defined and that a nonzero first Chern number obstructs a global
frame.
\end{exercise}

\subsection{Integer quantum Hall effect}

For a smooth gapped self-adjoint family $H(k)$ on the Brillouin torus, define
the occupied projection by
\[
P(k)=\frac{1}{2\pi i}\int_\Gamma(z-H(k))^{-1}\,dz.
\]
Its images form the \textbf{Bloch bundle} $E_{\mathrm{occ}}$.

\begin{exercise}[Gap-preserving Hall integer]
Prove that a gap-preserving homotopy produces isomorphic endpoint Bloch
bundles and preserves
\[
C_1(P)=\frac{i}{2\pi}\int_{\mathbb T^2}
\tr(P\,dP\wedge dP)\in\Z.
\]
Assuming the TKNN formula
$\sigma_{xy}=C_1(P)e^2/h$, compare the predicted plateaus with the integer
quantum Hall measurements \cite{tknn,von-klitzing}.  Identify gap closing as
the point where the invariant may change.
\end{exercise}

\subsection{Fredholm index}

A bounded operator $T:H_0\to H_1$ is \textbf{Fredholm} if its kernel and
cokernel are finite dimensional and its range is closed.  Define
\[
\operatorname{ind}T
=\dim\ker T-\dim\operatorname{coker}T.
\]

\begin{exercise}[Zero-mode imbalance is stable]
Prove that $T$ is Fredholm exactly when its image in the Calkin algebra is
invertible.  Deduce local constancy and invariance under compact perturbation
of the index.  State the Atiyah--Singer identification of analytic index with
the K-theoretic symbol pairing and derive the de Rham Euler-characteristic
case \cite{atiyah-singer}.
\end{exercise}

\begin{exercise}[Instanton charge produces chiral zero modes]
For an $SU(N)$ connection on a bundle over $S^4$, define
\[
Q(A)=\frac{1}{8\pi^2}\int_{S^4}\tr(F_A\wedge F_A).
\]
Use Chern--Weil theory to prove $Q(A)\in\Z$.  Assuming the coupled Dirac case
of the index theorem, relate $Q(A)$ to the difference of chiral zero-mode
dimensions and derive the integrated anomaly for $N_f$ massless fermions.
Distinguish the mathematical index identity from the physical assumptions in
the fermion model.
\end{exercise}


\section{Bell experiments and noncommutative analysis}

Why do entangled systems violate classical correlation bounds while remaining
unable to signal?

\subsection{Matrix norms and distinguishability}

A \textbf{concrete operator space} is a closed subspace
$E\subseteq\mathcal B(H,K)$ with matrix norms inherited from
\[
M_r(E)\subseteq\mathcal B(H^{\oplus r},K^{\oplus r}).
\]
For $\phi:E\to F$, define
\[
\|\phi\|_{\mathrm{cb}}=\sup_r\|1_{M_r}\otimes\phi\|.
\]

\begin{exercise}[Matrix levels detect noncommutativity]
Prove for transposition $t:M_n\to M_n$ that
\[
\|t\|=1,
\qquad
\|t\|_{\mathrm{cb}}=n.
\]
For the row and column spaces $R_n$ and $C_n$, prove that the identity between
their underlying Hilbert spaces has cb norm $\sqrt n$ in either direction.
Explain why an ordinary Banach norm forgets both distinctions.
\end{exercise}

For $u=\sum_kx_k\otimes y_k\in E\otimes F$, define the
\textbf{Haagerup norm}
\[
\|u\|_h
=\inf
\|[x_1\ \cdots\ x_r]\|_{M_{1,r}(E)}
\|[y_1\ \cdots\ y_r]^{\mathsf T}\|_{M_{r,1}(F)}.
\]

\begin{exercise}[Trace norm governs state discrimination]
Prove
\[
C_n\otimes_hR_n\cong M_n,
\qquad
R_n\otimes_hC_n\cong S_1^n.
\]
For equally likely states $\rho_0,\rho_1$, prove the Helstrom formula
\[
p_{\mathrm{succ}}
=\frac12+\frac14\|\rho_0-\rho_1\|_1.
\]
Interpret the trace norm through the second Haagerup product.
\end{exercise}

For $\Psi:M_n\to M_m$, define the \textbf{diamond norm}
\[
\|\Psi\|_\diamond
=\sup_r\|\Psi\otimes1_{M_r}\|_{1\to1}.
\]

\begin{exercise}[Entangled probes distinguish channels]
Prove stabilization at ancillary dimension $n$ and
\[
p_{\mathrm{succ}}(\Phi_0,\Phi_1)
=\frac12+\frac14\|\Phi_0-\Phi_1\|_\diamond.
\]
Identify the role of Schmidt rank and compare ancilla-assisted discrimination
with process-discrimination experiments \cite{watrous,sun}.
\end{exercise}

An \textbf{operator system} is a self-adjoint unital subspace of
$\mathcal B(H)$.  A unital map between operator systems is completely
positive when it preserves the positive cone at every matrix level.

\begin{exercise}[Operator-system positivity]
For an operator space $E\subseteq\mathcal B(H,K)$, form the Paulsen system
\[
\mathcal S(E)=
\left\{
\begin{pmatrix}\lambda1_K&x\\y^*&\mu1_H\end{pmatrix}:
x,y\in E,\ \lambda,\mu\in\C
\right\}.
\]
Prove that a map $\phi:E\to F$ is a complete contraction exactly when its
unital extension $\mathcal S(E)\to\mathcal S(F)$ is completely positive.
Explain why this gives a common language for channel admissibility and
matrix-level distinguishability.
\end{exercise}

\subsection{Bell correlations}

For self-adjoint unitaries $A_0,A_1,B_0,B_1$, define
\[
S=A_0\otimes(B_0+B_1)+A_1\otimes(B_0-B_1).
\]

\begin{exercise}[Quantum correlations reach the Tsirelson bound]
Prove that deterministic local sign assignments give $|S|\leq2$.  Using
\[
S^2=4-[A_0,A_1]\otimes[B_0,B_1],
\]
prove $\|S\|\leq2\sqrt2$ and attain equality with a maximally entangled
two-qubit state.  Compare the quantum prediction with loophole-controlled
Bell experiments \cite{tsirelson,hensen}.  Explain why violation excludes the
specified local hidden-variable model but does not permit signaling.
\end{exercise}

\begin{exercise}[Grothendieck bounds quantum advantage]
For real coefficients $c_{ij}$, compare
\[
\sup_{\varepsilon_i,\delta_j\in\{\pm1\}}
\left|\sum_{ij}c_{ij}\varepsilon_i\delta_j\right|
\quad\text{and}\quad
\sup_{\|x_i\|,\|y_j\|\leq1}
\left|\sum_{ij}c_{ij}\langle x_i,y_j\rangle\right|.
\]
Use Tsirelson's vector representation and the real Grothendieck inequality to
obtain a dimension-independent bound on bipartite correlation advantage.
Identify CHSH as a particular coefficient matrix.
\end{exercise}

\subsection{Infinite observables and index}

A \textbf{von Neumann algebra} is a unital $*$-subalgebra
$M\subseteq\mathcal B(H)$ equal to its bicommutant.  For a finite-factor
inclusion $N\subseteq M$, let $E_N:M\to N$ be the trace-preserving conditional
expectation and $e_N:L^2(M)\to L^2(N)$ the Jones projection.

\begin{exercise}[Weak limits produce the bicommutant]
Prove the finite-dimensional bicommutant theorem directly.  State the
von Neumann bicommutant theorem for a unital $*$-subalgebra of
$\mathcal B(H)$ and show that vector functionals are normal on its weak
closure.  Explain why weak closure, rather than norm closure alone, retains
limits of increasing families of quantum events.
\end{exercise}

\begin{exercise}[Subfactor index produces braid relations]
Prove $e_Nxe_N=E_N(x)e_N$.  In the basic-construction tower, derive
\[
e_i^2=e_i=e_i^*,
\qquad e_ie_{i\pm1}e_i=\delta^{-2}e_i,
\qquad e_ie_j=e_je_i\quad(|i-j|\ge2).
\]
Relate $\delta^2$ to the Jones index and derive braid generators from the
Temperley--Lieb relations.  State the quantization of subfactor indices below
four \cite{jones-index}.
\end{exercise}


\section{Observed anyonic braiding}

Why can exchanging experimentally realized quasiparticles in two dimensions
produce phases beyond bosonic and fermionic signs?

\subsection{Braids and fusion}

The \textbf{braid group} $B_n$ has generators
$\sigma_1,\ldots,\sigma_{n-1}$ and relations
\[
\sigma_i\sigma_{i+1}\sigma_i
=\sigma_{i+1}\sigma_i\sigma_{i+1},
\qquad
\sigma_i\sigma_j=\sigma_j\sigma_i\quad(|i-j|\ge2).
\]

A \textbf{$\C$-linear monoidal category} has a bilinear tensor product,
unit, and coherent associator.  A \textbf{braiding} is a natural family
$c_{x,y}:x\otimes y\to y\otimes x$ satisfying the hexagon identities.  In a
semisimple rigid category, decompositions
\[
x\otimes y\cong\bigoplus_zN_{xy}^zz
\]
are \textbf{fusion rules}; $\Hom(z,x\otimes y)$ is a fusion space.

\begin{exercise}[Braiding represents particle exchange]
Prove that
\[
\rho(\sigma_i)
=1^{\otimes(i-1)}\otimes c_{x,x}\otimes
1^{\otimes(n-i-1)}
\]
defines a representation of $B_n$ on $x^{\otimes n}$.  Identify the
Yang--Baxter equation with a hexagon consequence and prove that imposing
$\sigma_i^2=1$ gives the symmetric group.
\end{exercise}

\begin{exercise}[A three-cocycle controls reassociation]
For $G$-graded vector spaces, multiply the associator on degrees $g,h,k$ by
$\omega(g,h,k)\in U(1)$.  Prove that the pentagon identity is the
$3$-cocycle equation and that changing $\omega$ by a coboundary gives a
monoidally equivalent category.  For $G=\Z/2$, compute the two classes in
$H^3(G,U(1))$ and their pointed fusion rules.
\end{exercise}

The \textbf{monodromy} is $M_{x,y}=c_{y,x}c_{x,y}$.

\begin{exercise}[One-third anyons produce a $2\pi/3$ phase]
Assume all simple objects are invertible and form $\Z/m$, generated by $a$.
Prove that scalar monodromy compatible with fusion has the form
\[
M_{a^r,a^s}=\exp\!\left(\frac{2\pi ikrs}{m}\right).
\]
For $m=3$, $k=1$, derive the phase advance per enclosed anyon and compare it
with interferometric evidence for anyonic braiding \cite{nakamura}.
State the ideal topological assumptions and the non-topological device effects
that must be controlled.
\end{exercise}

\begin{exercise}[Fusion dimension]
Let positive numbers $d_x$ satisfy
$d_xd_y=\sum_zN_{xy}^zd_z$.  Prove that $d$ is a positive character of the
fusion ring.  For Fibonacci fusion
$\tau\otimes\tau=1\oplus\tau$, compute $d_\tau$ and the asymptotic growth of
the fusion space of $n$ $\tau$-anyons.  Relate nonintegral dimension to state
growth rather than to a fractional dimension of an ordinary vector space.
\end{exercise}

\subsection{Chern--Simons theory}

For a connection $A$ on an oriented three-manifold, define
\[
S_{\mathrm{CS}}(A)
=\frac{k}{4\pi}\int_M
\tr\!\left(A\wedge dA+\frac23A\wedge A\wedge A\right).
\]
A Wilson loop is $W_R(K)=\tr_R\operatorname{Hol}_A(K)$.

\begin{exercise}[Chern--Simons link invariants]
Vary $S_{\mathrm{CS}}$ and derive $F_A=0$.  Show that a large gauge
transformation changes the action by $2\pi k$ times an integer and deduce
level quantization.  Assuming the ribbon-category evaluation rules, prove that
closed braid amplitudes are categorical traces and recover the relation to
Jones-type invariants \cite{witten-jones,reshetikhin-turaev}.
\end{exercise}

\subsection{Stable functorial invariants}

An additive strong monoidal functor preserves direct sums and tensor products.
The \textbf{Grothendieck group} $K_0(\mathcal C)$ is generated by objects with
$[x\oplus y]=[x]+[y]$ and multiplication $[x][y]=[x\otimes y]$.

\begin{exercise}[Physical structures produce stable observables]
Prove that an additive strong monoidal functor induces a ring homomorphism on
$K_0$ and preserves categorical traces when it preserves duals.  Organize the
Aharonov--Bohm phase, Hall integer, Fredholm index, and anyon monodromy as
\[
\text{physical structure}
\longrightarrow\text{stable class}
\longrightarrow\text{pairing or trace}
\longrightarrow\text{observable}.
\]
For each, identify the experimental observable, the mathematical invariant,
and the hypothesis whose failure permits the value to change.
\end{exercise}

\begin{exercise}[Six experiments force six successive structures]
For each of the following observations, identify the minimal mathematical
object, the central invariant or theorem, and the model assumption tested by
the observation:
\begin{enumerate}
\item two-source interference and linear superposition;
\item water vibrational spectra and character multiplicities;
\item electromagnetic-wave propagation and $d^2=0$;
\item Aharonov--Bohm interference and flat holonomy;
\item integer Hall plateaus and the first Chern number;
\item Bell violation and matrix-level noncommutativity;
\item fractional-quantum-Hall interferometry and braid monodromy.
\end{enumerate}
For every step, exhibit a phenomenon that the preceding structure cannot
distinguish.  Prove that the resulting progression is
\[
\begin{gathered}
\text{linear spaces}\longrightarrow\text{representations}
\longrightarrow\text{geometry}\\[-2pt]
\longrightarrow\text{stable topology}\longrightarrow\text{operator spaces}
\longrightarrow\text{braided categories}.
\end{gathered}
\]
Separate observations from physical modeling assumptions and both from the
mathematical conclusions proved in the exercises.
\end{exercise}


\clearpage
\begin{thebibliography}{99}

\bibitem{stern-gerlach}
W. Gerlach and O. Stern,
\textit{Der experimentelle Nachweis der Richtungsquantelung im Magnetfeld},
Zeitschrift f"ur Physik 9 (1922), 349--352.

\bibitem{young}
T. Young,
\textit{The Bakerian Lecture: Experiments and calculations relative to
physical optics},
Philosophical Transactions of the Royal Society 94 (1804), 1--16.

\bibitem{choi}
M.-D. Choi,
\textit{Completely positive linear maps on complex matrices},
Linear Algebra and its Applications 10 (1975), 285--290.

\bibitem{stinespring}
W. F. Stinespring,
\textit{Positive functions on $C^*$-algebras},
Proceedings of the American Mathematical Society 6 (1955), 211--216.

\bibitem{altepeter}
J. B. Altepeter et al.,
\textit{Ancilla-assisted quantum process tomography},
Physical Review Letters 90 (2003), 193601.

\bibitem{wootters-zurek}
W. K. Wootters and W. H. Zurek,
\textit{A single quantum cannot be cloned},
Nature 299 (1982), 802--803.

\bibitem{shimanouchi}
T. Shimanouchi,
\textit{Tables of Molecular Vibrational Frequencies, Consolidated Volume I},
NSRDS--NBS 39, National Bureau of Standards, 1972.

\bibitem{hydrogen}
A. Kramida,
\textit{A critical compilation of experimental data on spectral lines and
energy levels of hydrogen, deuterium, and tritium},
Atomic Data and Nuclear Data Tables 96 (2010), 586--644.

\bibitem{hertz}
H. Hertz,
\textit{"Uber sehr schnelle elektrische Schwingungen},
Annalen der Physik 267 (1887), 421--448.

\bibitem{ligo}
B. P. Abbott et al.,
\textit{Observation of gravitational waves from a binary black hole merger},
Physical Review Letters 116 (2016), 061102.

\bibitem{aharonov-bohm}
Y. Aharonov and D. Bohm,
\textit{Significance of electromagnetic potentials in the quantum theory},
Physical Review 115 (1959), 485--491.

\bibitem{tonomura}
A. Tonomura et al.,
\textit{Evidence for Aharonov--Bohm effect with magnetic field completely
shielded from electron wave},
Physical Review Letters 56 (1986), 792--795.

\bibitem{berry-neutron}
T. Bitter and D. Dubbers,
\textit{Manifestation of Berry's topological phase in neutron spin rotation},
Physical Review Letters 59 (1987), 251--254.

\bibitem{tknn}
D. J. Thouless, M. Kohmoto, M. P. Nightingale, and M. den Nijs,
\textit{Quantized Hall conductance in a two-dimensional periodic potential},
Physical Review Letters 49 (1982), 405--408.

\bibitem{von-klitzing}
K. von Klitzing, G. Dorda, and M. Pepper,
\textit{New method for high-accuracy determination of the fine-structure
constant based on quantized Hall resistance},
Physical Review Letters 45 (1980), 494--497.

\bibitem{atiyah-singer}
M. F. Atiyah and I. M. Singer,
\textit{The index of elliptic operators I},
Annals of Mathematics 87 (1968), 484--530.

\bibitem{watrous}
J. Watrous,
\textit{Semidefinite programs for completely bounded norms},
Theory of Computing 5 (2009), 217--238.

\bibitem{sun}
K. Sun et al.,
\textit{Demonstration of Einstein--Podolsky--Rosen steering with enhanced
subchannel discrimination},
npj Quantum Information 4 (2018), 12.

\bibitem{tsirelson}
B. S. Tsirelson,
\textit{Quantum generalizations of Bell's inequality},
Letters in Mathematical Physics 4 (1980), 93--100.

\bibitem{hensen}
B. Hensen et al.,
\textit{Loophole-free Bell inequality violation using electron spins
separated by 1.3 kilometres},
Nature 526 (2015), 682--686.

\bibitem{jones-index}
V. F. R. Jones,
\textit{Index for subfactors},
Inventiones Mathematicae 72 (1983), 1--25.

\bibitem{nakamura}
J. Nakamura, S. Liang, G. C. Gardner, and M. J. Manfra,
\textit{Direct observation of anyonic braiding statistics},
Nature Physics 16 (2020), 931--936.

\bibitem{witten-jones}
E. Witten,
\textit{Quantum field theory and the Jones polynomial},
Communications in Mathematical Physics 121 (1989), 351--399.

\bibitem{reshetikhin-turaev}
N. Reshetikhin and V. G. Turaev,
\textit{Invariants of 3-manifolds via link polynomials and quantum groups},
Inventiones Mathematicae 103 (1991), 547--597.

\end{thebibliography}

\end{document}
